\section{Horizon Inference}
\label{sec:horizon_inference}

Inferential use requires more than the existence of a horizon. A data-driven
memory rule needs estimators of the drift parameters $(H,\zeta)$ together with a
criterion that compares inferential resolution to the curvature of the
useful-memory region.

Write $\alpha = \log \zeta$, and let $m$ denote the per-lag sample size used to
estimate lagged transport discrepancies. For a fixed lag design
$1,\dots,L$, define the lag-energy information scale
\[
  \mathcal I_{m,L}(H) \coloneqq m\sum_{j=1}^L j^{2H}.
\]

\begin{assumption}[Lag-geometry inference input]
\label{ass:lag_geometry_input}
For the lag model $D_j = \zeta j^H$, suppose there are estimators
$(\widehat \alpha_{m,L},\widehat H_{m,L})$ such that, on compact parameter sets,
\[
  \sqrt{\mathcal I_{m,L}(H)}
  \begin{pmatrix}
    \widehat \alpha_{m,L} - \alpha \\
    \widehat H_{m,L} - H
  \end{pmatrix}
  \Rightarrow \mathcal N(\mathbf 0, \Sigma_{L,H,\zeta}).
\]
\end{assumption}

The power-law lag-geometry model of \citet{Parreira2026Jul} provides one
concrete source of this assumption.

The continuous optimizer of the horizon envelope can be written as the smooth
map
\[
  \begin{aligned}
    n_\star(\alpha,H)
    &= \left(\frac{aC_K}{H C_S e^{\alpha}}\right)^{1/(a+H)}, \\
    g(\alpha,H) &\coloneqq \log n_\star(\alpha,H).
  \end{aligned}
\]
Its gradient is
\[
  \nabla g(\alpha,H)
  =
  \begin{pmatrix}
    -(a+H)^{-1} \\
    -\bigl[H(a+H)\bigr]^{-1} - g(\alpha,H)(a+H)^{-1}
  \end{pmatrix}.
\]

\begin{theorem}[Plug-in Horizon Central Limit Law]
\label{thm:plugin_horizon_clt}
Under \Cref{ass:lag_geometry_input},
\[
  \sqrt{\mathcal I_{m,L}(H)}
  \bigl(\log \widehat n_\star - \log n_\star\bigr)
  \Rightarrow \mathcal N(0,\tau^2),
\]
where $\widehat n_\star = n_\star(\widehat \alpha_{m,L},\widehat H_{m,L})$ and
\[
  \tau^2
  = \nabla g(\alpha,H)^\top \Sigma_{L,H,\zeta} \nabla g(\alpha,H).
\]
\end{theorem}

\begin{proof}
The map $g$ is smooth on any compact set bounded away from $H=0$ and
$\zeta=0$. Apply the multivariate delta method to
$(\widehat \alpha_{m,L},\widehat H_{m,L})$ under
\Cref{ass:lag_geometry_input}.
\end{proof}

The useful-memory region is the set
\[
  \mathcal U_\delta = \{n : \Phi(n) \le (1+\delta)\Phi(n_\star)\}.
\]
Let $(x_-(\delta;a,H),x_+(\delta;a,H))$ be the normalized endpoints from
\Cref{prop:useful_memory_region}, and define the log-radius
\[
  r_\delta(a,H)
  \coloneqq
  \min\bigl\{-\log x_-(\delta;a,H),\,\log x_+(\delta;a,H)\bigr\}.
\]
This quantity is the largest symmetric log-error around $n_\star$ that stays
inside the useful-memory region.

\begin{theorem}[Operational Identifiability of Useful Memory]
\label{thm:operational_identifiability}
Under \Cref{ass:lag_geometry_input},
\[
  \mathbb P\bigl(\widehat n_\star \in \mathcal U_\delta\bigr)
  =
  2\Phi_{\mathcal N}\!\left(
    \frac{\sqrt{\mathcal I_{m,L}(H)}\,r_\delta(a,H)}{\tau}
  \right) - 1 + o(1),
\]
where $\Phi_{\mathcal N}$ is the standard normal CDF. In particular,
\[
  \frac{\sqrt{\mathcal I_{m,L}(H)}\,r_\delta(a,H)}{\tau(H,\zeta)} \to \infty
  \quad \Longrightarrow \quad
  \mathbb P\bigl(\widehat n_\star \in \mathcal U_\delta\bigr) \to 1.
\]
\end{theorem}

\begin{proof}
By \Cref{prop:useful_memory_region},
$\widehat n_\star \in \mathcal U_\delta$ if and only if
$|\log \widehat n_\star - \log n_\star| \le r_\delta(a,H)$. Standardize this
event with \Cref{thm:plugin_horizon_clt} and use the symmetry of the Gaussian
limit.
\end{proof}

The quantity
\[
  \mathfrak S_{\delta,m,L}
  \coloneqq
  \frac{\sqrt{\mathcal I_{m,L}(H)}\,r_\delta(a,H)}{\tau(H,\zeta)}
\]
is therefore the operational identifiability score of the horizon. It compares
lag-geometry information to the curvature-imposed precision requirement of the
useful-memory region.

The following algorithm summarizes the operational pipeline.
\begin{algorithm}[!htbp]
\caption{Operational horizon-inference pipeline}
\label{alg:horizon_inference}
\footnotesize
\begin{algorithmic}[1]
\State Estimate recent lag discrepancies on lags $1,\dots,L$.
\State Fit the lag model to obtain $(\widehat\alpha_{m,L},\widehat H_{m,L})$ and a covariance proxy.
\State Compute $\widehat n_\star$ and $\widehat{\mathfrak S}_{\delta,m,L}$.
\If{$\widehat{\mathfrak S}_{\delta,m,L}$ is below the target tolerance threshold}
  \State widen to the useful-memory band or defer the update
\Else
  \State deploy $\widehat n_\star$ and monitor realized tracking error
\EndIf
\end{algorithmic}
\end{algorithm}

\begin{proposition}[Deterministic stability of the horizon map]
\label{prop:deterministic_horizon_stability}
Let $K\subset \mathbb R\times(0,\infty)$ be compact. Then there exists a finite
constant $L_K$ such that, for all $(\alpha_1,H_1),(\alpha_2,H_2)\in K$,
\[
  |g(\alpha_1,H_1)-g(\alpha_2,H_2)|
  \le L_K\bigl(|\alpha_1-\alpha_2|+|H_1-H_2|\bigr).
\]
Consequently, if $(\alpha,H)$ and $(\widehat\alpha_{m,L},\widehat H_{m,L})$ lie in
$K$ and
\[
  |\widehat\alpha_{m,L}-\alpha|+|\widehat H_{m,L}-H|
  \le \frac{r_\delta(a,H)}{L_K},
\]
then $\widehat n_\star\in \mathcal U_\delta$.
\end{proposition}

\begin{proof}
The gradient of $g$ is continuous on $K$, so $\|\nabla g\|_1$ is bounded there.
Apply the mean-value theorem to obtain the Lipschitz bound. The inclusion claim
follows from
$|\log \widehat n_\star-\log n_\star|\le r_\delta(a,H)$ together with the
characterization of $\mathcal U_\delta$ from \Cref{prop:useful_memory_region}.
\end{proof}

The informative regime lies across the transition from low to moderate
identifiability. Table \Cref{tab:inferable_horizon} shows representative
points across that transition.

\begin{table}[!htbp]
\centering
\small
\caption{Representative regimes across the inferable-horizon transition.}
\label{tab:inferable_horizon}
\begin{tabular}{rrrrrrr}
\toprule
$L$ & $m$ & $H$ & $\delta$ & $\mathfrak S$ & Empirical & Gaussian \\
\midrule
5 & 20 & 0.3 & 0.050 & 0.501 & 0.367 & 0.384 \\
8 & 10 & 0.9 & 0.020 & 1.001 & 0.664 & 0.683 \\
3 & 200 & 0.3 & 0.010 & 1.996 & 0.949 & 0.954 \\
8 & 100 & 0.3 & 0.005 & 3.001 & 0.997 & 0.997 \\
\bottomrule
\end{tabular}
\end{table}


On the tested grid, the plug-in central limit law is also well calibrated at the
scale level: the empirical standard deviation of $\log \widehat n_\star$ tracks
the delta-method prediction closely, and the nominal $95\%$ Gaussian interval
coverage stays near its target across the reported regimes. The simulation suite
also calibrates the regret expansion directly.

\begin{proposition}[Deterministic and asymptotic plug-in horizon regret]
\label{prop:plugin_horizon_regret}
Let $u = \log(\widehat n_\star/n_\star)$. Then, as $u \to 0$,
\[
  \frac{\Phi(\widehat n_\star)-\Phi(n_\star)}{\Phi(n_\star)}
  = \Psi(e^u)-1
  = \frac{H e^{-au} + a e^{Hu}}{a+H} - 1.
\]
Consequently, for every $r>0$,
\[
  |u|\le r
  \quad\Longrightarrow\quad
  \frac{\Phi(\widehat n_\star)-\Phi(n_\star)}{\Phi(n_\star)}
  \le
  B_r(a,H),
\]
where
\[
  B_r(a,H)
  \coloneqq
  \max\bigl\{\Psi(e^r)-1,\,\Psi(e^{-r})-1\bigr\}.
\]
Moreover,
\[
  \frac{\Phi(\widehat n_\star)-\Phi(n_\star)}{\Phi(n_\star)}
  = \frac{aH}{2}u^2 + O(u^3).
\]
Consequently,
\[
  \mathbb E\!
  \left[
    \frac{\Phi(\widehat n_\star)-\Phi(n_\star)}{\Phi(n_\star)}
  \right]
  = \frac{aH}{2}\frac{\tau^2}{\mathcal I_{m,L}(H)} + o\!\left(\mathcal I_{m,L}(H)^{-1}\right).
\]
\end{proposition}

\begin{proof}
By \Cref{prop:useful_memory_region},
\[
  \frac{\Phi(n_\star e^u)}{\Phi(n_\star)} = \Psi(e^u)
  = \frac{H e^{-au} + a e^{Hu}}{a+H}.
\]
The deterministic bound follows by maximizing this exact expression over
$u\in[-r,r]$. Differentiate at $u=0$. The first derivative vanishes and the
second derivative equals $aH$, giving the quadratic expansion. The expectation
bound then follows from \Cref{thm:plugin_horizon_clt} and the boundedness of
third moments on compact parameter sets.
\end{proof}

\begin{figure}[!htbp]
\centering
\safeincludegraphics[width=0.94\linewidth]{artifacts/figures/inferable_horizon/fig_operational_identifiability}
\caption{Operational identifiability of the useful-memory region. Binning by the operational score reveals a monotone rise in useful-region hit rate, with each $\delta$ curve tracking the Gaussian approximation from \Cref{thm:operational_identifiability}.}
\label{fig:inferable_horizon}
\end{figure}

Structural existence and operational identifiability need not coincide. The
relevant scale is the score $\mathfrak S_{\delta,m,L}$, which organizes the
transition from low coverage to near-certain useful-memory recovery.
